
We can now state the following theorem (\textup{[Per75]}, V, \SGARefLink{prop:sga3-V-3.1}{3.1}
and \SGARefLink{thm:sga3-VIA-3.2}{3.2}; see also \textup{[Per76]}, 0.0 and \SGARefLink{sga3:VIA:0.1}{0.1}).

\medskip
\phantomsection\label{thm:sga3-VIA-6.5}%
\noindent\textbf{Theorem 6.5 (D. Perrin).---}
\emph{Let \(G\) be a quasi-compact \(k\)-group.
\begin{enumerate}
\item[\textup{(i)}]\phantomsection\label{concl:sga3-VIA-6.5-i} \(G\) is the projective limit of a filtered
  system \((G_i)\) of \(k\)-groups of finite type (whose transition
  morphisms \(u_{ij}:G_j\to G_i\) are affine for \(i\) sufficiently
  large), and the morphisms \(G\to G_i\) are faithfully flat (and
  affine for \(i\) sufficiently large).
\item[\textup{(ii)}]\phantomsection\label{concl:sga3-VIA-6.5-ii} Let \(H\) be a closed \(k\)-subgroup of \(G\).
  Then the \(\mathrm{fpqc}\) quotient sheaf
  \(\widetilde{G/H}\) is a \(k\)-scheme in either of the following
  two cases:
  \begin{enumerate}
  \item[\textup{(1)}] the immersion \(H\to G\) is of finite
    presentation; in this case \(G/H\) is of finite type over \(k\);
  \item[\textup{(2)}] \(H\) is normal in \(G\).
  \end{enumerate}
\end{enumerate}}

For the proof of this theorem, which rests on several intermediate
theorems, we refer to \textup{[Per75]}. For the reader's
convenience, however, we prove the two corollaries below; cf.\
\textup{[Per75]}, V, \SGARefLink{par:sga3-VIA-3.3}{3.3}--3.4, or \textup{[Per76]}, 4.2.3--4.2.5.

\medskip
\phantomsection\label{cor:sga3-VIA-6.6}%
\noindent\textbf{Corollary 6.6.---}
\emph{Let \(f:G\to H\) be a quasi-compact morphism of \(k\)-groups.
\begin{enumerate}
\item[\textup{(i)}] If \(f\) is schematically dominant, then it is
  faithfully flat.
\item[\textup{(ii)}] This is so, in particular, if \(H\) is affine
  and the morphism
  \[
    f^\natural:\mathcal O(H)\longrightarrow\mathcal O(G)
  \]
  is injective.
\end{enumerate}}

\emph{Proof.}
Suppose that \(f\) is schematically dominant. By \hyperref[concl:sga3-VIA-6.2-i]{6.2(i)}, \(f\) is
surjective, so by \hyperref[lem:sga3-VIA-2.5.3]{2.5.3} it is enough to prove that \(f\) is flat at
the identity element of \(G\). Since
\[
  \mathcal O_{H,e}=\mathcal O_{H^0,e}
\]
(cf.\ \hyperref[thm:sga3-VIA-2.6.5]{2.6.5}), we may replace \(H\) by \(H^0\), and hence suppose
that \(H\) is irreducible. Then \(H\) is quasi-compact
(loc.\ cit.); since \(f\) is quasi-compact, \(G\) is also
quasi-compact. By \hyperref[concl:sga3-VIA-6.5-i]{6.5(i)},
\[
  H=\varprojlim_i H_i,
\]
where each \(H_i\) is an algebraic \(k\)-group. Let \(N_i\) be the
kernel of \(G\to H_i\); it is a closed normal \(k\)-subgroup of
\(G\). Moreover, since the identity section
\(\{e\}\to H_i\) is of finite presentation, the immersion
\(N_i\to G\) is of finite presentation. Hence, by \hyperref[concl:sga3-VIA-6.5-ii]{6.5(ii)}, the
\(\mathrm{fpqc}\) quotient \(G/N_i\) is an algebraic \(k\)-group.
We then have a commutative diagram
\begin{center}
  \phantomsection\label{diag:sga3-VIA-6.6-square}%
  \begin{tikzcd}[row sep=large,column sep=large]
    G
      \arrow[r,"f"]
      \arrow[d,"q_i"'] &
    H
      \arrow[d,"p_i"] \\
    G/N_i
      \arrow[r,"f_i"'] &
    H_i
  \end{tikzcd}
\end{center}
in which \(f\) is schematically dominant and \(p_i,q_i\) are
faithfully flat, hence schematically dominant. Thus
\[
  f_i\circ q_i=p_i\circ f
\]
is schematically dominant, and therefore so is \(f_i\). On the
other hand, \(f_i\) is a monomorphism, hence a closed immersion
because \(G/N_i\) and \(H_i\) are algebraic (\hyperref[prop:sga3-VIA-2.5.2]{2.5.2}). It follows that
\(f_i\) is an isomorphism, and consequently \(G\to H_i\) is
faithfully flat. By \textup{[BAC]}, I, \S~\SGARefLink{prop:sga3-VIB-2.7}{2.7}, Proposition~9, the
morphism \(G\to H\) is then flat; it is also surjective by \hyperref[concl:sga3-VIA-6.2-i]{6.2(i)},
and is therefore faithfully flat. This proves \SGARefLink{prop:sga3-VIB-2.7:item:i}{(i)}. Assertion \SGARefLink{prop:sga3-VIB-2.7:item:ii}{(ii)}
follows, because if \(H\) is affine and \(f^\natural\) is
injective, then the closed image of \(f\) is \(H\), so \(f\) is
schematically dominant (cf.\ \hyperref[prop:sga3-VIA-6.4]{6.4}).

\medskip
\phantomsection\label{cor:sga3-VIA-6.7}%
\noindent\textbf{Corollary 6.7.---}
\emph{Let \(u:G\to H\) be a morphism of \(k\)-groups and set
\(N=\ker(u)\). Suppose that \(u\) is quasi-compact.
\begin{enumerate}
\item[\textup{\SGAInlineRefTarget{cor:sga3-VIA-6.7:item:i}(i)}] The \(\mathrm{fpqc}\) quotient sheaf
  \(\widetilde{G/N}\) is represented by a \(k\)-group scheme
  \(G/N\), and \(u\) factors as
  \begin{center}
    \phantomsection\label{diag:sga3-VIA-6.7-factorization}%
    \begin{tikzcd}[ampersand replacement=\&,row sep=large,column sep=large]
      G \arrow[r,"u"] \arrow[d,"p"'] \& H \\
      G/N \arrow[ur,"i"']
    \end{tikzcd}
  \end{center}
  where \(p\) is faithfully flat and \(i\) is a closed immersion.
\item[\textup{\SGAInlineRefTarget{cor:sga3-VIA-6.7:item:ii}(ii)}] In particular, if \(u\) is a monomorphism,
  then \(u\) is a closed immersion; and if \(u\) is schematically
  dominant, then \(u\) is faithfully flat.
\end{enumerate}}

\emph{Proof.}
\textup{\SGARefLink{cor:sga3-VIA-6.7:item:i}{(i)}} By \hyperref[prop:sga3-VIA-6.4]{6.4} and \hyperref[cor:sga3-VIA-6.6]{6.6}, the closed image \(G'\) of \(u\) is a
closed group subscheme of \(H\), and
\[
  p:G\longrightarrow G'
\]
is faithfully flat and quasi-compact.%
\begingroup
\renewcommand{\thefootnote}{\fnsymbol{footnote}}%
\footnote[1]{\label{note:sga3-VIA-6.7-image-ambient}Translator's
note: the French source prints ``a closed group subscheme of
\(G\)'' at this point. Since \(u:G\to H\), its closed image lies in
\(H\), as the displayed factorization also shows.}%
\endgroup{}
Clearly \(\ker(p)=\ker(u)=N\). Thus, by \SGARefLink{sga3:iv:scholium:3:3:2:1}{Expos\'e~IV, 3.3.2.1} and
5.1.7, \(G'\) represents the \(\mathrm{fpqc}\) quotient sheaf
\(\widetilde{G/N}\).

\textup{\SGARefLink{cor:sga3-VIA-6.7:item:ii}{(ii)}} The second assertion is contained in \hyperref[cor:sga3-VIA-6.6]{6.6}; let us
prove the first. If \(u\) is a monomorphism, then so is \(p\).
Thus \(p\) is both a monomorphism and an effective epimorphism, and
is therefore an isomorphism (cf.\ \SGARefLink{sga3:iv:remark:1:14}{Expos\'e~IV, 1.14}). This proves
\hyperref[cor:sga3-VIA-6.7]{6.7}.

Let us finally record the following corollaries
(cf.\ \textup{[Per76]}, 4.2.6--4.2.8).

\medskip
\phantomsection\label{cor:sga3-VIA-6.8}%
\noindent\textbf{Corollary 6.8.---}
\emph{The category of commutative quasi-compact \(k\)-group schemes
is abelian.}

Taking \hyperref[cor:sga3-VIA-6.7]{6.7} into account, the proof is analogous to that of \hyperref[thm:sga3-VIA-5.4.2]{5.4.2}.

\medskip
\phantomsection\label{cor:sga3-VIA-6.9}%
\noindent\textbf{Corollary 6.9.---}
\emph{If \(\operatorname{char}(k)=0\), every \(k\)-group scheme
\(G\) is geometrically reduced.}

Indeed, if \(g\in G\) and \(K\) is an algebraic closure of
\(\kappa(g)\), then
\[
  \mathcal O_{G,e}\otimes_k K
    \simeq
  \mathcal O_{G_K,g_K}.
\]
It is therefore enough to prove that
\(\mathcal O_{G,e}=\mathcal O_{G^0,e}\) is geometrically reduced.
We are thus reduced to the case in which \(G\) is connected, hence
quasi-compact by \hyperref[thm:sga3-VIA-2.6.5]{2.6.5}. The result then follows from \hyperref[thm:sga3-VIA-6.5]{6.5} and
Cartier's theorem for algebraic groups
(cf.\ Expos\'e~VI\(_B\), \SGARefLink{cor:sga3-VIB-1.6.1}{1.6.1}, or \textup{[DG70]}, \S~II.6,
\SGARefLink{par:sga3-VIA-1.1}{Theorem~1.1}).

\medskip
\phantomsection\label{cor:sga3-VIA-6.10}%
\noindent\textbf{Corollary 6.10.---}
\emph{Let \(G\) be a quasi-compact \(k\)-group, and suppose that
\(k\) is algebraically closed.
\begin{enumerate}
\item[\textup{(i)}] If \(f:G\to H\) is a faithfully flat morphism
  of \(k\)-groups, then the induced map \(G(k)\to H(k)\) is
  surjective.
\item[\textup{(ii)}] The set of rational points is dense in \(G\).
\end{enumerate}}

For the proof, we refer to \textup{[Per75]}, V, \SGARefLink{prop:sga3-VIB-3.7}{3.7} and 3.9.

\section*{Bibliography}
\phantomsection\label{bib:sga3-VIA}
\addcontentsline{toc}{section}{Bibliography}
\footnote{\label{note:sga3-VIA-59}Editor's note: additional
references cited in this Expos\'e.}

\begin{description}
\item[{[An73]}] S. Anantharaman,
  \emph{Sch\'emas en groupes, espaces homog\`enes et espaces
  alg\'ebriques sur une base de dimension 1},
  M\'em. Soc. Math. France \textbf{33} (1973), 5--79.
\item[{[BAC]}] N. Bourbaki,
  \emph{Alg\`ebre commutative}, Chapters I--IV and V--VII,
  Masson, 1985.
\item[{[DG70]}] M. Demazure and P. Gabriel,
  \emph{Groupes alg\'ebriques}, Masson \& North-Holland, 1970.
\item[{[Per75]}] D. Perrin,
  \emph{Sch\'emas en groupes quasi-compacts sur un corps},
  Publ. Math. Orsay no.~165--75.46 (part 1),
  \texttt{http://portail.mathdoc.fr/PMO/}.
\item[{[Per76]}] D. Perrin,
  \emph{Approximation des sch\'emas en groupes, quasi-compacts
  sur un corps}, Bull. Soc. Math. France \textbf{104} (1976),
  323--335.
\item[{[Ray67a]}] M. Raynaud,
  \emph{Passage au quotient par une relation d'\'equivalence
  plate}, pp.~78--85 in \emph{Proc. Conf. Local Fields
  (Driebergen)}, ed.\ T.~A. Springer, Springer-Verlag, 1967.
\item[{[Ray67b]}] M. Raynaud,
  \emph{Sur le passage au quotient par un groupo\"\i de plat},
  C. R. Acad. Sci. Paris (S\'er. A) \textbf{265} (1967),
  384--387.
\item[{[Ray70]}] M. Raynaud,
  \emph{Faisceaux amples sur les sch\'emas en groupes et les
  espaces homog\`enes}, Lecture Notes in Mathematics
  \textbf{119}, Springer-Verlag, 1970.
\end{description}

